% =====================================================================
%  Chapter 5: Causal Inference and Identification
%  Source: P2 transcription (complete)
% =====================================================================

\chapter{Causal Inference and Identification}
\label{chap_Causal}

\section{The Goals of Empirical Research}
\label{sec_Goals}

Not all empirical questions are causal. \citet{Angrist2009a} organise the
goals of applied econometric work into a hierarchy:

\begin{figure}[htp]
\centering
\begin{tikzpicture}[>=Stealth,
    node distance=1.2cm and 2.2cm,
    every node/.style={draw, rounded corners, align=center, font=\small,
                       minimum width=2.8cm, minimum height=0.65cm}]
  % Root
  \node (root) {Empirical research};
  % Level 1
  \node (sum)  [below left=of root,  xshift=0.5cm]  {Summarising data};
  \node (inf)  [below right=of root, xshift=-0.5cm] {Inference};
  \draw[->] (root) -- (sum);
  \draw[->] (root) -- (inf);
  % Level 2 under Inference
  \node (desc) [below left=of inf,  xshift=0.3cm]  {Descriptive};
  \node (cf)   [below right=of inf, xshift=-0.3cm] {Counterfactual};
  \draw[->] (inf) -- (desc);
  \draw[->] (inf) -- (cf);
  % Level 3 under Counterfactual
  \node (pred) [below left=of cf,   xshift=0.5cm]  {Prediction\\(what-if)};
  \node (caus) [below right=of cf,  xshift=-0.5cm] {Causal\\identification};
  \draw[->] (cf) -- (pred);
  \draw[->] (cf) -- (caus);
\end{tikzpicture}
\caption{Taxonomy of the goals of empirical research.}
\label{fig_EmpiricalGoals}
\end{figure}

The distinction between descriptive inference and causal identification is
central. A descriptive regression summarises conditional means:
$\E[Y\mid X=x]$ tells us what values of $Y$ tend to accompany given values
of $X$, but nothing about what would happen to $Y$ if we \emph{changed} $X$.
Causal identification requires an additional assumption---a source of
exogenous variation that allows us to move from correlation to causation.

\subsection{Standard Regression Notation}

We can write the regression model in several equivalent ways:
\begin{alignat*}{2}
  Y_i &= m(\mathbf{x}) + \varepsilon && \\
      &= X_i\beta + \varepsilon_i && \quad\text{(standard form)} \\
  Y_i &\sim f_N(y_i \mid \mu_i, \sigma^2), \quad \mu_i = X_i\beta
    && \quad\text{(distributional form)}
\end{alignat*}

More generally, for a response variable $y_i$ following distribution
$f(y_i \mid \theta_i, \alpha)$ with systematic component $\theta_i =
g(X_i, \beta)$:

\begin{table}[htp]
\centering
\caption{Generalised regression notation across model families}
\label{tbl_RegNotation}
\renewcommand{\arraystretch}{1.4}
\begin{tabular}{@{}lllll@{}}
\toprule
\textbf{Model} & \textbf{Distribution} $f$ & \textbf{Link} $g$ &
\textbf{Range of $y$} & \textbf{Estimator} \\
\midrule
OLS          & Normal    & Identity  & $(-\infty, +\infty)$ & OLS / MLE \\
Logit        & Bernoulli & Logistic  & $\{0,1\}$             & MLE \\
Probit       & Bernoulli & Normal CDF & $\{0,1\}$            & MLE \\
Poisson      & Poisson   & Log       & $\{0,1,2,\ldots\}$   & MLE \\
Neg.\ Binomial & NB      & Log       & $\{0,1,2,\ldots\}$   & MLE \\
Gamma        & Gamma     & Log / Inv & $(0,+\infty)$         & MLE \\
Tobit        & Normal    & Identity  & $[0,+\infty)$         & MLE \\
Spatial Lag  & Normal    & Identity  & $(-\infty,+\infty)$   & MLE / IV \\
\bottomrule
\end{tabular}
\renewcommand{\arraystretch}{1}
\end{table}

\subsection{Two Types of Uncertainty}

\begin{itemize}
  \item \textbf{Estimation uncertainty:} arises because we observe a finite
    sample; standard errors shrink as $n \to \infty$.
  \item \textbf{Fundamental uncertainty:} arises because $Y$ is inherently
    stochastic; this uncertainty does not shrink with $n$. It is reflected in
    the residual variance $\sigma^2$.
\end{itemize}

\begin{figure}[htp]
\centering
\begin{tikzpicture}[>=Stealth, scale=1.0]
  % Three panels: linear, logistic, exponential systematic components
  % Panel 1: Linear
  \begin{scope}[xshift=0cm]
    \draw[->] (-0.2,0) -- (2.8,0) node[right] {\small$x$};
    \draw[->] (0,-0.2) -- (0,2.4) node[above] {\small$\mu$};
    \draw[very thick] (0,0.3) -- (2.5,2.2);
    \node[below] at (1.25,-0.3) {\small Linear};
  \end{scope}
  % Panel 2: Logistic
  \begin{scope}[xshift=4.0cm]
    \draw[->] (-0.2,0) -- (2.8,0) node[right] {\small$x$};
    \draw[->] (0,-0.2) -- (0,2.4) node[above] {\small$\mu$};
    \draw[very thick, domain=0:2.5, samples=60]
      plot (\x, {2.1/(1+exp(-4*(\x-1.25)))});
    \draw[dashed, gray] (-0.1, 2.1) -- (2.7, 2.1);
    \node[below] at (1.25,-0.3) {\small Logistic};
  \end{scope}
  % Panel 3: Exponential
  \begin{scope}[xshift=8.0cm]
    \draw[->] (-0.2,0) -- (2.8,0) node[right] {\small$x$};
    \draw[->] (0,-0.2) -- (0,2.4) node[above] {\small$\mu$};
    \draw[very thick, domain=0:2.3, samples=60]
      plot (\x, {0.3*exp(1.0*\x)});
    \node[below] at (1.25,-0.3) {\small Exponential};
  \end{scope}
\end{tikzpicture}
\caption{Three systematic component forms: linear (identity link), logistic
(logit link), and exponential (log link).}
\label{fig_SystematicComp}
\end{figure}


\section{The Fundamental Problem of Causal Inference}
\label{sec_FPCI}

\subsection{Potential Outcomes and the Rubin Causal Model}

Let $D_i \in \{0,1\}$ be a binary treatment indicator. Each unit $i$ has
two \textbf{potential outcomes}\index{potential outcomes}:
\begin{itemize}
  \item $Y_{0i}$: outcome if untreated ($D_i = 0$)
  \item $Y_{1i}$: outcome if treated ($D_i = 1$)
\end{itemize}

The observed outcome is (\citealt{Holland1986}; Rubin Causal Model):

\begin{tcolorbox}[keyresult]
\begin{equation}
  Y_i = Y_{0i} + (Y_{1i} - Y_{0i})D_i
  \label{eq_PotentialOutcomes}
\end{equation}
\end{tcolorbox}

The \textbf{individual treatment effect} is $Y_{1i} - Y_{0i}$. The
\textbf{fundamental problem of causal inference} is that we observe only one
potential outcome for each unit---either $Y_{0i}$ or $Y_{1i}$, never both.

\subsection{The Average Treatment Effect on the Treated (ATT)}

\begin{tcolorbox}[keyresult]
\begin{align}
  \underbrace{\E[Y_i \mid D_i=1] - \E[Y_i \mid D_i=0]}_{\text{observed difference}}
  &= \underbrace{\E[Y_{1i} - Y_{0i} \mid D_i=1]}_{\text{ATT}}
   + \underbrace{\E[Y_{0i} \mid D_i=1] - \E[Y_{0i} \mid D_i=0]}_{\text{selection bias}}
  \label{eq_ATT}
\end{align}
\end{tcolorbox}

The observed difference between treated and untreated units confounds the
causal effect (ATT) with selection bias---the difference in potential
untreated outcomes between those who were and were not treated.


\section{The Selection Problem in Urban Contexts}
\label{sec_Selection}

\subsection{The Hospital Paradox}

A striking illustration: patients in hospitals have worse health outcomes than
patients not in hospitals. Does this mean hospitals are harmful? No---sicker
people are more likely to \emph{select into} hospital treatment. The
selection mechanism determines who receives treatment, creating the
confounding. The same logic applies throughout urban and regional analysis:
neighbourhoods, cities, and regions are not random assignments.

\begin{table}[htp]
\centering
\caption{The selection problem: conventional estimates vs.\ corrected estimates}
\label{tbl_Selection1}
\renewcommand{\arraystretch}{1.4}
\begin{tabular}{@{}lcc@{}}
\toprule
& \textbf{Treated (hospital)} & \textbf{Untreated (no hospital)} \\
\midrule
Average health outcome       & Low    & High \\
Na\"{i}ve treatment effect   & \multicolumn{2}{c}{Negative (spurious)} \\
True treatment effect (ATT)  & \multicolumn{2}{c}{Positive (causal)} \\
Selection bias               & \multicolumn{2}{c}{Negative (sicker select in)} \\
\bottomrule
\end{tabular}
\renewcommand{\arraystretch}{1}
\end{table}

\begin{table}[htp]
\centering
\caption{REO (real estate owned) vs.\ conventional property sales:
a selection problem in housing markets}
\label{tbl_Selection2}
\renewcommand{\arraystretch}{1.4}
\begin{tabular}{@{}lcc@{}}
\toprule
Property characteristic & REO sale & Conventional sale \\
\midrule
Seller motivation       & Distressed (bank) & Voluntary \\
Condition at sale       & Often degraded    & Market-maintained \\
Observed sale price     & Discounted        & Full market value \\
Na\"{i}ve ``REO discount'' & Negative      & Baseline \\
True REO effect         & \multicolumn{2}{c}{Requires identification strategy} \\
\bottomrule
\end{tabular}
\renewcommand{\arraystretch}{1}
\end{table}

\subsection{Geographic Sorting and the Endogeneity of Place}

Cities and neighbourhoods are not random assignments---they are the outcomes
of sorting by households, firms, and governments over decades and centuries.
When we observe that neighbourhood $A$ has better outcomes than neighbourhood
$B$, we cannot distinguish the causal effect of neighbourhood characteristics
from the selection process that determined who lives where. The standard
spatial regression coefficient conflates both. This is the core identification
challenge of regional policy informatics.


\section{The Experimental Ideal}
\label{sec_Experimental}

\subsection{Random Assignment}

Under random assignment, $D_i$ is independent of potential outcomes:
\begin{equation}
  (Y_{0i}, Y_{1i}) \ind D_i
\end{equation}
This implies $\E[Y_{0i}\mid D_i=1] = \E[Y_{0i}\mid D_i=0]$, so the
selection bias term in equation~\eqref{eq_ATT} vanishes:
\begin{equation}
  \E[Y_i \mid D_i=1] - \E[Y_i \mid D_i=0] = \text{ATT}
  = \E[Y_{1i} - Y_{0i} \mid D_i=1]
\end{equation}

Note: randomisation \emph{spreads} selection bias symmetrically across
treatment and control groups---it does not remove selection bias from
individual units but makes it zero \emph{in expectation}.

\paragraph{Example: The Perry Pre-School Programme (Heckman et al., 2010)}
\index{Perry Pre-School}

The Perry Pre-School Programme (1962) was a randomised controlled trial
(RCT) that randomly assigned disadvantaged African-American children in
Ypsilanti, Michigan to either a high-quality pre-school programme (treatment)
or a control group. Follow-up studies by \citet{Heckman2010} traced
participants into adulthood, finding substantial long-run effects on education,
employment, criminal behaviour, and earnings---effects that far exceeded the
short-run IQ gains that were the original target of the programme. The Perry
RCT remains one of the most influential examples of a field experiment in
social science, demonstrating both the power of random assignment and the
importance of measuring long-run outcomes.

\subsection{Regression with Constant Treatment Effect}

Under random assignment and assuming constant treatment effects:

\begin{tcolorbox}[keyresult]
\begin{equation}
  Y_i = \alpha + \beta D_i + \eta_i
  \label{eq_TreatmentReg}
\end{equation}
where $\beta = \E[Y_{1i} - Y_{0i}]$ is the Average Treatment Effect (ATE)
and $\alpha = \E[Y_{0i}]$ is the mean counterfactual outcome.
\end{tcolorbox}


\section{The Conditional Independence Assumption}
\label{sec_CIA}

When randomisation is not available, causal identification typically requires
one of three strategies: CIA (this section), instrumental variables
(Chapter~\ref{chap_BeyondOLS}), or quasi-experimental designs
(Chapter~\ref{chap_BeyondOLS}).

\begin{tcolorbox}[keyresult]
\textbf{Conditional Independence Assumption (CIA):}
\begin{equation}
  \{Y_{0i}, Y_{1i}\} \ind D_i \mid X_i
  \label{eq_CIA}
\end{equation}
Selection into treatment is independent of potential outcomes \emph{after
conditioning on observable covariates} $X_i$.
\end{tcolorbox}

Under the CIA, the regression with control variables $X_i$:
\begin{equation}
  Y_i = \alpha + \beta D_i + X_i\gamma' + \eta_i
\end{equation}
\emph{has a causal interpretation}---$\beta$ estimates the average treatment
effect conditional on $X_i$.

\noindent\textbf{Limitations of the CIA.} The CIA requires that we have
observed and included all confounders---variables that affect both treatment
assignment and outcomes. In the presence of \emph{unobserved} confounders,
the CIA fails and regression estimates are biased. This is the case in most
spatial settings, where geographic sorting creates selection on unobservables
(amenity preferences, ability, social capital) that are not captured in
administrative data.


\section{The TSS Decomposition and $R^2$}
\label{sec_TSS}

The decomposition of total variation into explained and residual components
follows directly from the orthogonality of $\hat{\mathbf{y}}$ and
$\hat{\mathbf{e}}$:

\begin{align*}
  \sum(y_i - \bar{y})^2
  &= \sum(\hat{y}_i - \bar{y})^2 + \sum(y_i - \hat{y}_i)^2 \\
  \text{TSS} &= \text{ESS} + \text{RSS}
\end{align*}

\noindent\textit{Proof:}
\begin{align*}
  \sum(y_i - \bar{y})^2
  &= \sum\bigl[(\hat{y}_i-\bar{y}) + \hat{e}_i\bigr]^2 \\
  &= \sum(\hat{y}_i-\bar{y})^2 + 2\sum(\hat{y}_i-\bar{y})\hat{e}_i
     + \sum\hat{e}_i^2 \\
  &= \sum(\hat{y}_i-\bar{y})^2 + \sum\hat{e}_i^2
\end{align*}
where the cross term vanishes because $\hat{\mathbf{y}}'\hat{\mathbf{e}} = 0$
and $\boldsymbol{\iota}'\hat{\mathbf{e}} = 0$ (the constant is always
included, so the residuals sum to zero). $\square$

\renewcommand\bibname{Further reading}
\begin{subappendices}
\bibliographystyle{econometrica}
\bibliography{Literature/OverLord}
\IfFileExists{Literature/Causal.bbl}{\input{Literature/Causal.bbl}}{\typeout{Need bibtex on Causal}}
\end{subappendices}
